\section{Introduksjon}
\label{sec:introduksjon}
\subsection{Bakgrunn og formål}
Prosjektet omfatter utvikling av en mekanisk lekebil med svinghjulsbasert energilagring. Arbeidet høsten 2026 skal resultere i en fungerende fysisk prototype og en rapport som dokumenterer systematisk produktutvikling, 3D-design, produksjon og grunnleggende testing, med et omfang tilsvarende semesteroppgaven i MAST1101.

\subsection{Problemstilling}
\begin{quote}
Hvordan kan en mekanisk lekebil med svinghjulsbasert energilagring utvikles gjennom systematisk produktutvikling, 3D-design og prototyping for å oppnå størst mulig kjøredistanse ved ett standardisert opptrekk?
\end{quote}
Målet er å utvikle og sammenligne praktiske løsninger innenfor prosjektets rammer. Gjennomsnittlig kjøredistanse er hovedmålet, supplert med repeterbarhet, driftssikkerhet og rettlinjet kjøring.

\subsection{Avgrensning og leveranser høsten 2026}
Dette semesteret omfatter hele bilen: kravspesifikasjon, konseptutvikling og konseptvalg, CAD-modellering, produksjon, montering og grunnleggende eksperimentell testing. Rapporten utvikles samtidig med prototypen, slik at designvalg, endringer og forsøk blir dokumentert fortløpende.

Detaljert dynamisk og kinematisk modellering, beregning av tidsavhengig bevegelse og modellbasert optimalisering av drivverket avgrenses til våren 2027. Høstens tester skal evaluere funksjon og ytelse og gi et sammenligningsgrunnlag for videre arbeid.
